\section{Muitos requests, memória compartilhada, vários núcleos}
\label{sec:muitos-requests}

\paragraph{O problema.}
Num servidor de servlets, várias threads executam o mesmo código. A
\citetitle{doc:jakarta-servlet61}, §2.3.3.1, diz que
\enquote{A servlet container may send concurrent requests through the
service method of the servlet}, e o desenvolvedor do servlet precisa
prever a execução de várias threads nesse método. A mesma seção
recomenda não sincronizar o método \texttt{service}, por causa da
perda de desempenho: \enquote{It is strongly recommended that
Developers not synchronize the service method}.

O texto \citetitle{artigo:kegel-c10k}, de Dan Kegel, traz copyright de
1999 a 2018 e última atualização em 5 de fevereiro de 2014. A abertura pergunta: \enquote{It's
time for web servers to handle ten thousand clients simultaneously,
don't you think?} O texto defende que atender dez mil clientes ao mesmo
tempo pede algo além de uma thread por cliente, e diz sobre o Java:
\enquote{Up through JDK 1.3, Java's standard networking libraries
mostly offered the one-thread-per-client model.}

\paragraph{A decisão.}
A \citetitle{doc:jsr51}, com Final Release em 9 de maio de 2002, entrou
no Java~1.4 (2002). O líder da JSR foi Mark Reinhold, e Kegel e Doug
Lea estavam no grupo de especialistas. A descrição da JSR lista
\enquote{APIs for scalable I/O, fast buffered binary and character I/O,
regular expressions, charset conversion, and an improved filesystem
interface}. Segundo Kegel, o \texttt{java.nio} do JDK~1.4 trouxe ao
Java o I/O não bloqueante.

\paragraph{O que mudou depois.}
Depois veio o padrão reativo. O \citetitle{doc:reactive-streams} se
define como \enquote{a standard for
asynchronous stream processing with non-blocking back pressure}. A
\termo{back pressure}{Contrapressão: o consumidor diz ao produtor quanto
consegue receber, para não ser inundado.} é parte da definição. A
versão~1.0.0 para a JVM
(\citetitle{doc:reactive-streams-jvm}) é de 29 de abril de 2015. O
JDK~9 trouxe a interface \texttt{java.util.concurrent.Flow}
(\enquote{Since: 9}), cujo javadoc diz: \enquote{These interfaces
correspond to the reactive-streams specification}
(\citetitle{doc:javase9-flow}).

O \citetitle{doc:jep444} descreve o custo desse caminho. No estilo
assíncrono, o código devolve a thread a um pool enquanto espera I/O, e
o fim da operação é avisado por um \termo{callback}{Função passada
para outra, que a chama quando um evento acontece, como o fim de uma
operação de I/O.}. A lógica de cada requisição vira uma sequência de
estágios pequenos, em geral lambdas, compostos por uma API como
\texttt{CompletableFuture} ou um framework reativo, e o programador
abre mão de construções como laços e blocos \texttt{try}/\texttt{catch}
(\enquote{They thus forsake the language's basic sequential composition
operators}). Um dos objetivos do JEP é que servidores no estilo
thread-per-request escalem com virtual threads, finais no JDK~21
(seção~\ref{sec:threads-ate-vt}).

No Quarkus, o guia
\citetitle{doc:quarkus-reactive-architecture} diz que \enquote{A few
threads can handle many concurrent I/O}. Uma requisição HTTP é
processada numa dessas poucas threads de I/O, que por isso não podem
bloquear (\enquote{you can't block the thread anymore}). Código
bloqueante é desviado para \termo{worker threads}{Threads de um pool
separado, usadas para executar código que bloqueia sem ocupar as
threads de I/O.}: \enquote{Quarkus implements a proactor pattern that
switches to worker thread when needed}. O guia
\citetitle{doc:quarkus-virtual-threads} diz que \enquote{the support of
virtual thread is implemented using the @RunOnVirtualThread
annotation}; o que a anotação faz está na
seção~\ref{sec:virtual-threads}.

\paragraph{O que isso muda no seu código.}
Cada recomendação tem seu escopo.
\begin{itemize}
\item \emph{Jakarta Servlet~6.1, cujos containers exigem Java SE~17 ou
  mais novo (§1.2):} o mesmo servlet pode atender
  requisições concorrentes pelo método \texttt{service}, então escreva
  esse código prevendo várias threads (seções~\ref{sec:monitores}
  e~\ref{sec:jmm-historia}). A especificação recomenda não
  sincronizar o método \texttt{service}.
\item \emph{JDK~9 em diante:} \texttt{java.util.concurrent.Flow} é a
  versão do JDK das interfaces do Reactive Streams.
\item \emph{JDK~21 em diante:} código bloqueante no estilo
  thread-per-request pode escalar com virtual threads, sem reescrever
  a lógica em estágios assíncronos (seção~\ref{sec:threads-ate-vt}).
\item \emph{Quarkus~3, sobre JDK~17 ou mais novo
  (\citetitle{doc:quarkus-getting-started}):} código nas threads de I/O não pode bloquear, e
  código bloqueante vai para as worker threads; com Java~21 ou mais
  novo, \texttt{@RunOnVirtualThread} é a terceira opção
  (seção~\ref{sec:vt-vs-reativo}).
\end{itemize}

\paragraph{Onde aprofundar.}
Quando escolher virtual threads, reativo ou worker threads no Quarkus
é o tema da seção~\ref{sec:vt-vs-reativo}. As regras de uso das
virtual threads estão na seção~\ref{sec:virtual-threads}.
